% Generated draft: requires mathematical and source review.
\documentclass[11pt]{article}
\usepackage[margin=25mm]{geometry}
\usepackage{fontspec}
\setmainfont{Latin Modern Roman}
\usepackage{amsmath,amssymb,mathtools,unicode-math}
\setmathfont{Latin Modern Math}
\usepackage{enumitem,needspace,xurl,hyperref}
\setlength{\parindent}{0pt}
\setlength{\parskip}{5pt}
\setlength{\emergencystretch}{3em}
\newcommand{\N}{\mathbb{N}}
\newcommand{\Z}{\mathbb{Z}}
\newcommand{\Q}{\mathbb{Q}}
\newcommand{\R}{\mathbb{R}}
\newcommand{\C}{\mathbb{C}}
\newcommand{\F}{\mathbb{F}}
\newcommand{\A}{\mathbb{A}}
\newcommand{\PP}{\mathbb{P}}
\newcommand{\E}{\mathbb{E}}
\newcommand{\eps}{\varepsilon}
\newcommand{\dd}{\,\mathrm{d}}
\begin{document}
\section*{UnsolvedMath ID: 30000814}
\textbf{Source problem code:} catalogue entry
\subsection{Definitions and mathematical statement}
Let \(X\) be a smooth projective toric variety of dimension \(d\), let \(A\) be ample and \(B\) nef, and retain the intended adjoint divisor \(D\). The supplied clean reconstruction proposes
\[
 D_{\mathrm{cand}}=(d+2)A+B+K_X,
\]
whereas the original/statement fields print \(D=A^{d+2}+B+K_X\). The verification explicitly says the OCR expression could not be recovered unambiguously and assigns reconstructed-unverified status. Tensor-power notation for a line bundle could explain an additive multiple, but that interpretation is not certified by the record. The candidate is therefore conditional, not promoted to a recovered divisor formula.

For a specified integral torus-invariant Cartier representative \(D=\sum_\rho a_\rho D_\rho\), with character lattice \(M\) and primitive ray generators \(v_\rho\), its polytope convention is
\[
 P_D=\{u\in M\otimes_\Z\R:\langle u,v_\rho\rangle\geq-a_\rho\text{ for all }\rho\}.
\]
The source presupposes an associated lattice polytope but does not give this representative or confirm the divisor expression. For a lattice polytope \(P\), normality here means the integer-decomposition property
\[
 mP\cap M=\underbrace{(P\cap M)+\cdots+(P\cap M)}_{m\text{ summands}}
 \quad\text{for every integer }m\geq1.
\]

A unimodular covering means a finite collection of lattice simplices \(\Delta_i\subseteq P\) whose union is \(P\), each unimodular in the relevant lattice. For a full-dimensional \(d\)-simplex with vertices \(u_0,\ldots,u_d\), unimodular means that \(u_1-u_0,\ldots,u_d-u_0\) form a \(\Z\)-basis of \(M\), equivalently normalized volume one. If \(P\) has smaller dimension, the affine-lattice convention must be stated. A covering need not be a face-to-face triangulation, regular triangulation or disjoint-interior decomposition, and no such stronger conclusion is substituted.

For the intended adjoint family, assuming \(P_D\) normal, the question is whether \(P_D\) has a unimodular covering. It does not ask whether every arbitrary normal lattice polytope does. The original attributes normality to a cited vanishing theorem, while the clean version takes it as a hypothesis; the theorem and its assumptions are not reproduced, so normality is not independently asserted.

The 2026-08-21 literature assessment says the divisor/polytope was omitted and no general answer located. The current reconstruction supplies a tentative expression but the verification still withholds confirmation. The open label is retained as dated metadata, not a verified current theorem status.
\subsection{Short English statement}
For the source's intended normal adjoint toric polytope, ask whether it can be covered by unimodular lattice simplices. A covering is weaker than a triangulation. The divisor expression is explicitly reconstructed and unverified, so the cleaned multiple \((d+2)A\) remains a candidate rather than a silently accepted OCR repair.
\subsection{Sources}
\begin{itemize}
\item Local UnsolvedMath registry, permanent ID 30000814, source code catalogue entry, ``Unimodular Coverings of Adjoint Toric Polytopes,'' including the original \(A^{d+2}+B+K_X\) expression, candidate cleaned formula and reconstructed-unverified warning.
\item \emph{Mini-Workshop: Projective Normality of Smooth Toric Varieties (2008)}, Oberwolfach Reports, pp.\ 2283--2320, DOI 10.4171/owr/2007/39: \url{https://doi.org/10.4171/owr/2007/39}. This is the sole supplied literature reference. The citation uses 2008, while the DOI and proposed-year metadata use 2007.
\item The normality claim refers to a vanishing theorem not identified in the supplied record; no separate citation is invented.
\item Literature triage checked 2026-08-21; record updated 2026-09-16. No external source has been visited or independently checked here.
\end{itemize}
\end{document}
